% بخش: روش‌ها
% فصل: اثبات‌ها
% بخش فرعی: مقدمه

\documentclass[../../../include/open-logic-section]{subfiles}

\begin{document}

\olfileid{mth}{prf}{int}

\olsection{مقدمه}

با توجه به تجربه‌تان در منطق مقدماتی، شاید با دستگاهی از نوع
!!a{derivation} راحت باشید — احتمالاً دستگاهِ !!{derivation} طبیعی
یا دستگاهی به سبک فیچ، و شاید هم دستگاه درخت اثبات. احتمالاً به یاد
دارید که در این دستگاه‌ها اثبات انجام می‌دادید: یا !!a{formula} را
اثبات می‌کردید یا نشان می‌دادید استدلالی معین معتبر است. برای این
کار قواعد دستگاه را به کار می‌بردید تا به نتیجهٔ نهاییِ مطلوب
برسید. هنگام استدلال \emph{دربارهٔ} منطق نیز چیزهایی را اثبات
می‌کنیم، اما در بیشتر موارد از دستگاهی مبتنی بر !!a{derivation}
استفاده نمی‌کنیم. در واقع، بیشتر اثبات‌هایی که بررسی می‌کنیم به زبان
فارسی انجام می‌شوند (گاه با اندکی زبان نمادی در میان)، نه سراسر به
زبان منطق مرتبه‌اول. هنگام ساختن چنین اثبات‌هایی شاید در آغاز
سردرگم شوید: بدون دستگاهی از نوع !!a{derivation} چگونه چیزی را
اثبات کنم؟ از کجا آغاز کنم؟ چگونه بدانم اثباتم درست است؟

پیش از کوشش برای یک اثبات، مهم است بدانیم اثبات چیست و چگونه ساخته
می‌شود. همان‌گونه که نامش نشان می‌دهد، \emph{اثبات} باید نشان دهد
چیزی صادق است. می‌توانید این را همچون گفت‌وگویی در نظر بگیرید: کسی
از شما می‌پرسد آیا گزاره‌ای صادق است؛ مثلاً آیا هر عدد اولی جز دو
فرد است. پاسخِ صرفِ «بله» کافی نیست؛ شاید بخواهد بداند \emph{چرا}.
در این صورت، برای او اثباتی عرضه می‌کنید.

در گفتار روزمره شاید اشاره‌ای به پاسخ یا پاسخی ناتمام کافی باشد.
اما در منطق و ریاضیات اثباتی دقیق می‌خواهیم؛ می‌خواهیم صدق چیزی را
فراتر از \emph{هرگونه} تردید نشان دهیم. این بدان معناست که هر گام
اثبات باید موجه باشد و توجیه نیز استوار باشد (یعنی فرضی که به کار
می‌برید واقعاً در صورت قضیهٔ مورد اثبات فرض شده باشد، تعریف‌ها
درست به کار روند، توجیه‌های مورد استناد استنتاج‌های درستی باشند، و
جز این‌ها).

معمولاً گزاره‌ای را اثبات می‌کنیم. گزاره‌های مورد اثبات را با
نام‌های گوناگون می‌خوانیم: حکم، قضیه، لم یا نتیجه. حکم گزاره‌ای پایه
و شایستهٔ اثبات است: آن‌قدر مهم هست که ثبت شود، اما شاید چندان ژرف
نباشد یا چندان به کار نرود. قضیه حکمی مهم و معنادار است. اثبات آن
اغلب به چند گام شکسته می‌شود و گاهی نام کسی را می‌گیرد که نخست آن
را اثبات کرده است (برای نمونه، قضیهٔ کانتور یا قضیهٔ
L\"owenheim--Skolem)، یا نام واقعیتی را که بدان مربوط است (برای
نمونه، قضیهٔ تمامیت). لم حکم یا قضیه‌ای است که در اثبات نتیجه‌ای
مهم‌تر به کار می‌رود. گیج‌کننده اینکه گاهی لم‌ها خود نتایجی مهم‌اند
و نام کسی را می‌گیرند که آن‌ها را معرفی کرده است (برای نمونه، لم
زورن). نتیجه، حکمی است که به‌آسانی از حکمی دیگر به دست می‌آید.

گزاره‌ای که باید اثبات شود اغلب فرض‌هایی دارد که روشن می‌کنند
دربارهٔ چه نوع چیزهایی مطلبی را اثبات می‌کنیم. شاید با «فرض کنید
$!A$ یک !!a{formula} به صورت $!B \lif !C$ باشد» یا «فرض کنید
$\Gamma \Proves !A$» یا عبارتی مانند این‌ها آغاز شود. این‌ها
\emph{فرض‌های} حکم، قضیه یا لم هستند و می‌توانید در اثبات خود
صادقشان فرض کنید. آن‌ها دامنهٔ ادعای ما را محدود می‌کنند و
نام‌هایی برای اشیای مورد بحث نیز معرفی می‌کنند. برای نمونه، اگر حکم
شما با «فرض کنید $!A$ یک !!a{formula} به صورت
$!B \lif !C$ باشد» آغاز شود، تنها دربارهٔ همهٔ فرمول‌های نوعی خاص
(یعنی شرطی‌ها) چیزی را اثبات می‌کنید، و چنین فهمیده می‌شود که
$!B \lif !C$ یک شرطیِ دلخواه است که اثباتتان دربارهٔ آن سخن خواهد
گفت.

\end{document}
